% !TeX root = ../main.tex
\chapter{Detalhamento do Problema}
\label{cap:problema}

Historicamente, o código-fonte é produzido por desenvolvedores humanos,
que escrevem manualmente cada linha, aplicando suas habilidades e
conhecimentos para construir software funcional. Esse processo envolve
a compreensão dos requisitos do sistema, a implementação de algoritmos e
estruturas de dados e a realização de testes para garantir o
funcionamento correto do software. Mesmo com desenvolvedores experientes,
entretanto, erros podem ocorrer e resultar em vulnerabilidades de
segurança.

Nos últimos anos, esse cenário tem sido transformado pelo advento de
ferramentas de Inteligência Artificial voltadas ao desenvolvimento, como
GitHub Copilot, Claude Code, Codex e Cursor, que utilizam modelos de
linguagem avançados para auxiliar os desenvolvedores na escrita de
código. Essas ferramentas podem gerar código automaticamente — desde
pequenos trechos até projetos inteiros construídos por agentes a partir
de descrições em linguagem natural. Embora aumentem a produtividade e
reduzam o tempo de desenvolvimento, também introduzem novos desafios de
segurança.

Estudos empíricos e levantamentos da indústria apontam que o código gerado ou assistido por IA pode apresentar um volume expressivo de falhas de segurança --- com relatórios recentes indicando taxas de até 4,4 vezes mais vulnerabilidades em cenários específicos de desenvolvimento acelerado por assistentes inteligentes ---, com determinados tipos de fraquezas aparecendo com frequência e \textit{templates} inseguros se repetindo em diferentes repositórios, o que sugere padrões induzidos pelos próprios modelos \cite{stateofcyber2026,wu2025}. 

Contudo, o desafio da segurança de código transcende a origem do artefato: seja o código escrito puramente por humanos ou produzido com auxílio de modelos gerativos, a necessidade de validação antes do \textit{merge} e da integração definitiva à \textit{branch} principal é imperativa. As ferramentas de análise estática tradicionalmente utilizadas para identificar esses problemas apresentam duas limitações centrais: elevada taxa de falsos positivos e dificuldade em compreender o contexto do código, como a intenção do desenvolvedor e o fluxo de dados entre funções \cite{yang2025,techscience2025}. Diante disso, há demanda por abordagens mais contextuais, baseadas na integração de análise estática a classificadores de \textit{Machine Learning} e, de forma complementar, \textit{Large Language Models}, capazes de detectar vulnerabilidades e priorizar alertas com maior precisão no fluxo contínuo de desenvolvimento \cite{wang2025,mdpi2026}.

